\subsubsection{Data Persistence}
\begin{enumerate}
    \item \textbf{R Interface for Array Serialization and Deserialization}
          
          The R interface provides bindings to the Fortran serialization and deserialization
          routines, enabling efficient handling of multi-dimensional arrays. This interface
          supports R arrays, matrices, and vectors with automatic type conversion to
          compatible Fortran data types.
          \begin{enumerate}
            \item {General Architecture}
                  
                  The R package provides ten main functions for array handling:
                  
                  \begin{itemize}
                    \item \texttt{tox\_serialize\_int\_array(array, filename)} -
                          Serialize N-dimensional integer arrays
                          
                    \item \texttt{tox\_deserialize\_int\_array(filename)} - Deserialize
                          N-dimensional integer arrays
                          
                    \item \texttt{tox\_serialize\_real\_array(array, filename)} -
                          Serialize N-dimensional real arrays
                          
                    \item \texttt{tox\_deserialize\_real\_array(filename)} - Deserialize
                          N-dimensional real arrays
                          
                    \item \texttt{tox\_serialize\_char\_array(array, filename)} -
                          Serialize N-dimensional character arrays
                          
                    \item \texttt{tox\_deserialize\_char\_array(filename)} - Deserialize
                          N-dimensional character arrays
                          
                    \item \texttt{tox\_serialize\_logical\_array(array, filename)} -
                          Serialize N-dimensional logical array
                          
                    \item \texttt{tox\_deserialize\_logical\_array(filename)} -
                          Deserialize N-dimensional logical array
                          
                    \item \texttt{tox\_serialize\_complex\_array(array, filename)} -
                          Serialize N-dimensional complex array
                          
                    \item \texttt{tox\_deserialize\_complex\_array(filename)} -
                          Deserialize N-dimensional complex array
                  \end{itemize}
                  
            \item \textbf{Array Functions}
                  
                  \textbf{Serialization Functions:}
                  \begin{itemize}
                    \item \texttt{array} (input): R array, matrix, or vector to be serialized
                          \\ \textbf{Type:}
                          \begin{itemize}
                            \item \texttt{Integer}
                                  
                            \item \texttt{Numeric} (double precision)
                                  
                            \item \texttt{Characters}
                                  
                            \item \texttt{Logical}
                                  
                            \item \texttt{Complex}
                          \end{itemize}
                          \textbf{Dimensions:} Theoretically are N dimensions supported, but
                          typical use cases should not exceed 5 dimensions \\
                          
                    \item \texttt{filename} (input): Name of the binary file to create \\
                          \textbf{Type:} \texttt{character}
                  \end{itemize}
                  \begin{lstlisting}[language=R, caption={R Array Serialization Example}]
tox_serialize_int_array(array, filename)
tox_serialize_real_array(array, filename)
tox_serialize_char_array(array, filename)
                  \end{lstlisting}
                  
                  \textbf{Deserialization Functions:}
                  \begin{itemize}
                    \item \texttt{filename} (input): Name of the binary file to read \\ \textbf{Type:}
                          \texttt{character}
                          
                    \item \textbf{Return Value:} Deserialized R array with original
                          dimensions \\ \textbf{Type:} \texttt{array}, \texttt{matrix}, or \texttt{vector}
                          \\ \textbf{Data Type:} Same as original
                  \end{itemize}
                  \begin{lstlisting}[language=R, caption={R Array Deserialization Example}]
restored_int <- tox_deserialize_int_array(filename)
restored_real <- tox_deserialize_real_array(filename)
restored_char <- tox_deserialize_char_array(filename)
                  \end{lstlisting}
                  
            \item \textbf{Implementation Details}
                  
                  \begin{itemize}
                    \item The R interface automatically converts between R data types and
                          Fortran-compatible formats
                          
                    \item Integer vectors are converted to 32-bit integers
                          
                    \item Numeric vectors are converted to 64-bit floating point
                          
                    \item Character arrays handle variable-length strings with automatic
                          padding
                          
                    \item All array dimensions are preserved during serialization/deserialization
                          
                    \item The interface handles both regular arrays and matrices seamlessly
                  \end{itemize}
                  
                  \subsubsection*{Usage Example}
                  
                  \begin{lstlisting}[language=R, caption={Full example of serialization and deserialization in R}]
tox_serialize_int_array(arr1, fn("int1d.bin"))
restored <- tox_deserialize_int_array(fn("int1d.bin"))

# Real array example
arr2 <- matrix(runif(12), nrow = 3, ncol = 4)
tox_serialize_real_array(arr2, fn("real2d.bin"))
restored <- tox_deserialize_real_array(fn("real2d.bin"))

# Character array example
arr3 <- c("GENE1", "BRCA1", "TP53", "EGFR")
tox_serialize_char_array(arr3, fn("char1d.bin"))
restored <- tox_deserialize_char_array(fn("char1d.bin"))
                  \end{lstlisting}
                  
            \item \textbf{Compatibility Notes}
                  
                  \begin{itemize}
                    \item Requires R with appropriate Fortran bindings
                          
                    \item Supports standard R data types (integer, numeric, character, logical,
                          complex)
                          
                    \item Preserves array attributes and dimensions during serialization
                          
                    \item Character arrays maintain content but may have padding removed
                          upon deserialization
                  \end{itemize}
          \end{enumerate}
    
    \item \textbf{Create Zip Archive (Low-Level) (R)}

The low-level function creates a zip archive from keys and filenames,
directly calling the Fortran backend. Use this function for custom data serialization
scenarios or when you need fine-grained control over archive creation.
Call the \texttt{create\_zip\_archive} function with the following arguments:

    \begin{itemize}
    \item \texttt{zip\_filename}: Name of the zip file to create
          
    \item \texttt{keys}: Vector of keys for the manifest (identifiers for
          each file)
          
    \item \texttt{filenames}: Vector of filenames to include in archive
    \end{itemize}

\begin{lstlisting}[language=R, caption={R Create Zip Archive (Low-Level)}]
create_zip_archive(zip_filename, keys, filenames)
\end{lstlisting}

\item \textbf{Extract Zip Archive (R)}

The function extracts a zip archive created by \texttt{create\_zip\_archive}
and returns a list mapping data keys to extracted filenames. This is useful
for accessing individual files from custom archives. Call the \texttt{extract\_zip\_archive}
function with the following arguments:

\begin{itemize}
    \item \texttt{zip\_filename}: Path to the zip file to extract
\end{itemize}

\textbf{Returns:}
\begin{itemize}
    \item List mapping data keys to extracted temporary filenames
\end{itemize}

\begin{lstlisting}[language=R, caption={R Extract Zip Archive}]
extract_zip_archive(zip_filename)
\end{lstlisting}

\item \textbf{Save TOX Data (High-Level) (R)}

The high-level function saves TOX data to a zip archive, handling
validation, serialization, and archive creation for standard-conform tox
gene data. This function automatically manages temporary files and cleanup.
Call the \texttt{save\_tox\_data} function with the following arguments:

\begin{itemize}
    \item \texttt{zip\_filename}: Name of the zip file to create
          
    \item \textbf{(Optional)} \texttt{gene\_ids}: Character vector of gene
          IDs
          
    \item \textbf{(Optional)} \texttt{gene\_ids\_name}: Filename for gene
          IDs in the archive
          
    \item \textbf{(Optional)} \texttt{expression\_vectors}: Numeric matrix
          of expression values (n\_samples × n\_genes)
          
    \item \textbf{(Optional)} \texttt{expression\_vectors\_name}: Filename
          for expression vectors in the archive
          
    \item \textbf{(Optional)} \texttt{gene\_to\_fam}: Integer vector mapping
          each gene to its family index (0 = unassigned)
          
    \item \textbf{(Optional)} \texttt{gene\_to\_fam\_name}: Filename for
          gene to family mapping in the archive
          
    \item \textbf{(Optional)} \texttt{family\_ids}: Character vector of
          family IDs
          
    \item \textbf{(Optional)} \texttt{family\_ids\_name}: Filename for
          family IDs in the archive
          
    \item \textbf{(Optional)} \texttt{family\_centroids}: Numeric matrix of
          family centroids (n\_samples × n\_families)
          
    \item \textbf{(Optional)} \texttt{family\_centroids\_name}: Filename for
          family centroids in the archive
          
    \item \textbf{(Optional)} \texttt{shift\_vectors}: Numeric matrix of
          shift vectors (2*n\_samples × n\_genes)
          
    \item \textbf{(Optional)} \texttt{shift\_vectors\_name}: Filename for
          shift vectors in the archive
\end{itemize}

\begin{lstlisting}[language=R, caption={R Save TOX Data}]
save_tox_data(zip_filename,
              gene_ids = NULL, gene_ids_name = NULL,
              expression_vectors = NULL, expression_vectors_name = NULL,
              gene_to_fam = NULL, gene_to_fam_name = NULL,
              family_ids = NULL, family_ids_name = NULL,
              family_centroids = NULL, family_centroids_name = NULL,
              shift_vectors = NULL, shift_vectors_name = NULL)
\end{lstlisting}

\item \textbf{Read TOX Data (R)}

The function reads data from a zip archive created by \texttt{save\_tox\_data},
extracting and deserializing the requested data components for standard-conform
tox gene data. This function automatically handles file extraction and
cleanup. Call the \texttt{read\_tox\_data} function with the following
arguments:

\begin{itemize}
    \item \texttt{zip\_filename}: Name of the zip file to read from
          
    \item \textbf{(Optional)} \texttt{gene\_ids}: If not NULL, will attempt to
          read gene IDs from archive
          
    \item \textbf{(Optional)} \texttt{expression\_vectors}: If not NULL, will
          attempt to read expression vectors from archive
          
    \item \textbf{(Optional)} \texttt{gene\_to\_fam}: If not NULL, will attempt
          to read gene to family mapping from archive
          
    \item \textbf{(Optional)} \texttt{family\_ids}: If not NULL, will attempt
          to read family IDs from archive
          
    \item \textbf{(Optional)} \texttt{family\_centroids}: If not NULL, will attempt
          to read family centroids from archive
          
    \item \textbf{(Optional)} \texttt{shift\_vectors}: If not NULL, will attempt
          to read shift vectors from archive
\end{itemize}

\textbf{Return list:}
\begin{itemize}
    \item \textbf{gene\_ids}: Loaded gene IDs vector (if requested)
          
    \item \textbf{expression\_vectors}: Loaded expression vectors matrix (if
          requested)
          
    \item \textbf{gene\_to\_fam}: Loaded gene-to-family mapping vector (if
          requested)
          
    \item \textbf{family\_ids}: Loaded family IDs vector (if requested)
          
    \item \textbf{family\_centroids}: Loaded family centroids matrix (if
          requested)
          
    \item \textbf{shift\_vectors}: Loaded shift vectors matrix (if requested)
\end{itemize}

\begin{lstlisting}[language=R, caption={R Read TOX Data}]
read_tox_data(zip_filename,
              gene_ids = NULL,
              expression_vectors = NULL,
              gene_to_fam = NULL,
              family_ids = NULL,
              family_centroids = NULL,
              shift_vectors = NULL)
\end{lstlisting}

\textbf{R Data Archiving Workflow Examples:} 
\begin{lstlisting}[language=R, caption={R Standard TOX Data Archiving}]
# Save complete TOX dataset
save_tox_data(
    zip_filename = "complete_tox_dataset.zip",
    gene_ids = gene_ids,
    gene_ids_name = "gene_ids.bin",
    expression_vectors = expression_vectors,
    expression_vectors_name = "expression.bin",
    gene_to_fam = gene_to_fam,
    gene_to_fam_name = "gene_to_fam.bin",
    family_ids = family_ids,
    family_ids_name = "family_ids.bin",
    family_centroids = family_centroids,
    family_centroids_name = "centroids.bin",
    shift_vectors = shift_vectors,
    shift_vectors_name = "shift_vectors.bin"
)

# Save partial dataset
save_tox_data(
    zip_filename = "basic_data.zip",
    gene_ids = gene_ids,
    gene_ids_name = "gene_ids.bin",
    expression_vectors = expression_vectors,
    expression_vectors_name = "expression.bin"
)

# Read data back with selective loading
loaded_data <- read_tox_data(
    zip_filename = "complete_tox_dataset.zip",
    gene_ids = TRUE,
    expression_vectors = TRUE,
    gene_to_fam = TRUE
)
\end{lstlisting}

\begin{lstlisting}[language=R, caption={R Custom Archive Creation and Extraction}]
# Create custom archive with non-standard data
create_zip_archive(
    zip_filename = "custom_data.zip",
    keys = c("3d_tensor", "results"),
    filenames = c("tensor_data.bin", "analysis_results.bin")
)

# Extract and access custom archive
file_mapping <- extract_zip_archive("custom_data.zip")
cat("Files in archive:\n")
print(file_mapping)
# Output: $3d_tensor = "tensor_data.bin", $results = "analysis_results.bin"
\end{lstlisting}

\begin{lstlisting}[language=R, caption={R Mixed Standard and Custom Data Archiving}]
# Serialize custom arrays
tox_serialize_real_array(custom_array, "custom_data.bin")
tox_serialize_int_array(another_array, "another_data.bin")

# Create mixed archive with both standard and custom data
create_zip_archive(
    zip_filename = "mixed_archive.zip",
    keys = c("standard_expression", "custom_feature", 
        "analysis_results"),
    filenames = c("expression.bin", "custom_data.bin", 
        "another_data.bin")
)

# Later, extract and process
mapping <- extract_zip_archive("mixed_archive.zip")
expression_data <- tox_deserialize_real_array(
    mapping[["standard_expression"]])
custom_features <- tox_deserialize_real_array(
    mapping[["custom_feature"]])
results <- tox_deserialize_int_array(
    mapping[["analysis_results"]])
\end{lstlisting}

\textbf{Important Notes for R Archive Functions:}
\begin{itemize}
    \item All archive functions automatically handle temporary file cleanup
          
    \item The \texttt{save\_tox\_data} function validates data dimensions
          before serialization
          
    \item Both \texttt{save\_tox\_data} and \texttt{read\_tox\_data} use the
          standard TOX data keys: "gene\_ids", "expression", "gene\_to\_family",
          "family\_ids", "family\_centroids", "shift\_vectors"
          
    \item Archives created with R functions are compatible with Python and Fortran
          implementations
          
    \item Error codes are automatically checked and converted to informative
          error messages
\end{itemize}
\end{enumerate}